\chapter{Connection-Relative Informational Pullback Geometry}
\label{ch:pullback-geometry}

The preceding chapters separate the principal bundle, the two associated law
bundles, and the normalized probabilistic model.  This chapter adds a smooth
information-geometric layer without identifying those objects.  A chosen
connection converts the derivative of an agent section into a vertical first
jet.  The Fisher metric and the Amari--Chentsov tensor can then be pulled back
to the contextual base.  The resulting tensors are invariant under passive
changes of frame, but they depend on the selected connection and may be
degenerate.  These qualifications are part of the construction rather than
defects to be removed by notation.

The same construction also supplies the differential interface to
coarse-graining.  A coarse morphism whose scale-connection defect form takes
values in the isotropy subalgebra of the coarse section value intertwines the
covariant first jets, while a parameter-independent Markov channel contributes
a positive-semidefinite vertical Fisher defect.  That isotropy condition is the
only meaning this chapter attaches to compatibility of the two connections, and
it is stated as \eqref{eq:pb-isotropy-criterion} rather than left to the
reader.  The vertical defect composes by an exact unconditional cocycle
identity, whereas its pullback to the base composes only up to an exact
residual that this chapter displays.  Deterministic energy precomposition and
Galerkin aggregation do not satisfy the Markov hypothesis merely because they
are called coarse maps.

\section{Statistical tensors on the associated fibers}
\label{sec:pb-statistical-tensors}

\hypothesisheading{Regular channel models}{hyp:pb-regular-models}
For each channel $x\in\{b,m\}$, let $\mathcal B_x$ be a finite-dimensional
smooth parametrized-measure model satisfying the selected smooth-tier
hypothesis of \Cref{hyp:geo-smooth-tier}.  Assume differentiability in
quadratic mean, a positive-definite Fisher form, third-power integrability of
every score direction used below, smooth dependence of the resulting second-
and third-order tensors, and domination sufficient to differentiate the
divergence expressions below through third order.  Assume also that the represented action
$\widehat\rho_x(g)$ is induced by a parameter-independent bimeasurable change
of sample coordinates and preserves $\mathcal B_x$.  These assumptions exclude
singular statistical strata, changing-support models without the stated
differentiability, and arbitrary actions on parameter coordinates.
\status{HYPOTHESIS}

At $p\in\mathcal B_x$, write $\ell_u$ for the score in direction
$u\in T_p\mathcal B_x$.  Define
\begin{align}
g^F_{x,p}(u,v)
&=\E_p[\ell_u\ell_v],
\label{eq:pb-fiber-fisher}\\
\mathcal T_{x,p}(u,v,w)
&=\E_p[\ell_u\ell_v\ell_w].
\label{eq:pb-fiber-amari}
\end{align}
The first is the Fisher metric and the second is the symmetric
Amari--Chentsov tensor.  Their existence here is a statement about the
selected parametrized-measure models, not about arbitrary law fibers.
\status{DEFINITION}

\propositionheading{Descent of the statistical tensors}{prop:pb-statistical-tensor-descent}
Under \Cref{hyp:pb-regular-models}, the represented $G$-action
preserves $g_x^F$ and $\mathcal T_x$.  Consequently they descend to smooth
vertical tensors, denoted by the same symbols, on
$\mathcal E_x=P\times_{\widehat\rho_x}\mathcal B_x$.  \status{ESTABLISHED}

\paragraph{Proof.}
Let $r_g$ be the bimeasurable sample-coordinate change representing $g$.
Pushforward defines the unitary map
$L^2(p)\to L^2((r_g)_\#p)$ given by
$f\mapsto f\circ r_g^{-1}$.  Because $r_g$ is independent of the statistical
parameter, differentiability in quadratic mean identifies the score of the
pushed-forward tangent direction with $\ell_u\circ r_g^{-1}$.  The
pushforward integration identity preserves the second and third score moments in
\eqref{eq:pb-fiber-fisher} and \eqref{eq:pb-fiber-amari}.  An invariant tensor
on the fiber descends through the associated-bundle quotient.  Explicitly, let
$\iota_u:\mathcal B_x\to\mathcal E_x$ send $\beta$ to the class $[u,\beta]$ of
a principal frame $u$.  The quotient convention
\eqref{eq:geo-quotient-convention} gives
$\iota_{ug}=\iota_u\circ\widehat\rho_x(g)$, so a fiber tensor $\tau$ has
frame-independent transport $\tau^u:=(\iota_u^{-1})^*\tau$ exactly when
$\widehat\rho_x(g)^*\tau=\tau$ for every $g$.  Smoothness of the descended
tensor follows by evaluating this construction on local smooth sections of
$P$.  $\square$

The bimeasurable sample action in this proposition is load bearing.  Closure
of a set of parameter values under an arbitrary diffeomorphism of its
parameter chart would not prove statistical isometry and would not define the
same descended tensors.  \status{ESTABLISHED}

\section{Covariant first jets and pullback tensors}
\label{sec:pb-covariant-jets}

Fix one channel $x$ temporarily and write $E=\mathcal E_x$,
$\mathcal B=\mathcal B_x$, $\varpi=\varpi_x$, $g^F=g_x^F$,
$\mathcal T=\mathcal T_x$, and $\omega=\omega_x$.  The connection induces a
horizontal subbundle $H^\omega E\subset TE$ and a vertical projection
$\operatorname{ver}^{\omega}:TE\to VE$.  For a smooth section
$s:\mathcal U\to E$ over an open set $\mathcal U\subseteq\mathcal C$, define
\begin{equation}
D^\omega s
:=\operatorname{ver}^{\omega}\circ Ts
:T\mathcal U\longrightarrow s^*VE.
\label{eq:pb-covariant-first-jet}
\end{equation}
The connection-split first jet is the pair
$j^1_\omega s=(s,D^\omega s)$.  For nonlinear statistical fibers,
$D^\omega s$ is a vertical-valued one-form along $s$; it is not a linear
covariant derivative on a vector space of sections.  \status{DEFINITION}

\definitionheading{Covariant informational pullbacks}{def:pb-informational-pullbacks}
The connection-relative Fisher and Amari pullbacks of $s$ are
\begin{align}
h_s^\omega(X,Y)
&=g^F_{s(c)}(D^\omega s(X),D^\omega s(Y)),
\label{eq:pb-fisher-pullback}\\
c_s^\omega(X,Y,Z)
&=\mathcal T_{s(c)}(D^\omega s(X),D^\omega s(Y),D^\omega s(Z)),
\label{eq:pb-amari-pullback}
\end{align}
for $X,Y,Z\in T_c\mathcal C$.  Thus
$h_s^\omega\in\Gamma(\operatorname{Sym}^2T^*\mathcal U)$ and
$c_s^\omega\in\Gamma(\operatorname{Sym}^3T^*\mathcal U)$.  The first is called a
semimetric until its nondegeneracy is proved.  \status{DEFINITION}

\theoremheading{Passive gauge covariance of the covariant pullbacks}{thm:pb-pullback-gauge-invariance}
Under a passive local change of principal frame, with the connection and
section coordinates transformed by their associated laws, both
$h_s^\omega$ and $c_s^\omega$ are unchanged as tensors on $\mathcal U$.
\status{ESTABLISHED}

\paragraph{Proof.}
In a local frame, the vertical jet transforms by the tangent representation,
\begin{equation}
D^{\omega'}s'(X)
=T\widehat\rho_x(g^{-1})D^\omega s(X),
\label{eq:pb-covariant-jet-gauge-law}
\end{equation}
where $g:\mathcal U\to G$ is the frame change and the inverse follows from
the quotient convention \eqref{eq:geo-quotient-convention}.  The connection
law is \eqref{eq:geo-local-connection-b} or
\eqref{eq:geo-local-connection-m}, according to the channel.  By
\Cref{prop:pb-statistical-tensor-descent}, the tangent representation is an
isometry of $g^F$ and preserves $\mathcal T$.  Substitution in
\eqref{eq:pb-fisher-pullback} and \eqref{eq:pb-amari-pullback} cancels the
frame change in every argument.  $\square$

The word invariance in \Cref{thm:pb-pullback-gauge-invariance} means passive
covariance: two coordinate descriptions of the same pair $(\omega,s)$ give the
same tensor.  It does not say that two different connections give the same
pullback, and it does not say that the tensors are fixed by the active gauge
group at a fixed connection.  For the active statement, take the trivial
principal $\R$-bundle over $\R$ with the normal location fiber, the zero
connection, and the zero section, so that $h=0$; the bundle automorphism
$F(u(x)g)=u(x)(x+g)$ carries that section to $s'(x)=\mathcal N(x,1)$ and hence
$h$ to $dx^2$.  \status{ESTABLISHED}

\section{Exact dependence on the chosen connection}
\label{sec:pb-connection-dependence}

Let $\omega'$ be a second principal connection on $P$.  Their difference is
horizontal and $\operatorname{Ad}$-equivariant; denote its descended form by
\begin{equation}
a\in\Omega^1(\mathcal C,\operatorname{Ad}(P)).
\label{eq:pb-connection-difference}
\end{equation}
For the associated action, let
$\vartheta_e:\operatorname{Ad}(P)_{\varpi(e)}\to V_eE$ denote the fundamental
vertical map.  Along $s$, set
\begin{equation}
R_a^s(X)=\vartheta_{s(c)}(a_c(X)).
\label{eq:pb-connection-difference-vertical}
\end{equation}
The sign convention is fixed frame-free by
\eqref{eq:pb-connection-difference-vertical}: in a local trivialization the
horizontal lift of $X$ at $(c,\beta)$ is $(X,-\zeta_{A(X)}\beta)$, where
$\zeta_\xi$ is the fundamental vertical field of $\xi\in\mathfrak g$ for the
represented action, so replacing $A$ by $A+a$ adds $\vartheta(a(X))$ to the
vertical jet.  For a general law fiber $\mathcal B\subseteq\mathcal P(\mathsf K)$
the fiber is not a vector space and the linear expression
$\widehat\rho_{x*}(A)$ is undefined; it is available only in a declared
fiberwise-linear associated-vector-bundle realization, where it agrees with
$\zeta_{A(X)}$.  The convention is compatible with
\eqref{eq:geo-quotient-convention}.  \status{DEFINITION}

\propositionheading{Connection-change formulas}{prop:pb-pullback-connection-change}
The two vertical jets and their Fisher pullbacks obey
\begin{align}
D^{\omega'}s
&=D^\omega s+R_a^s,
\label{eq:pb-jet-connection-change}\\
h_s^{\omega'}(X,Y)
&=h_s^\omega(X,Y)
+g^F(R_a^sX,D^\omega sY)
+g^F(D^\omega sX,R_a^sY)
+g^F(R_a^sX,R_a^sY).
\label{eq:pb-fisher-connection-change}
\end{align}
The Amari pullback satisfies the exact polarized identity
\begin{equation}
c_s^{\omega'}(X,Y,Z)
=\mathcal T((D^\omega s+R_a^s)X,
(D^\omega s+R_a^s)Y,
(D^\omega s+R_a^s)Z).
\label{eq:pb-amari-connection-change}
\end{equation}
It therefore contains the original term, three terms with one occurrence of
$R_a^s$, three terms with two occurrences, and one term with three
occurrences.  \status{ESTABLISHED}

\paragraph{Proof.}
In one local frame, replacing the horizontal lift
$(X,-\zeta_{A(X)}\beta)$ by $(X,-\zeta_{A(X)+a(X)}\beta)$ shifts the vertical
part of $Ts$ by $\vartheta_{s(c)}(a(X))=R_a^s(X)$, which proves
\eqref{eq:pb-jet-connection-change}.  Bilinearity of $g^F$ gives
\eqref{eq:pb-fisher-connection-change}, and trilinearity of
$\mathcal T$ gives \eqref{eq:pb-amari-connection-change}.  The formulas are
independent of the selected local frame because $a$ is an
$\operatorname{Ad}(P)$-valued one-form and every displayed contraction is
gauge invariant.  $\square$

The dependence can occur even in the simplest regular family.  Let
$\mathcal C=\R$, let $P$ be the trivial principal translation bundle, and let
$\mathcal B=\{\mathcal N(\mu,1):\mu\in\R\}$.  For the constant section
$\mu=0$, the zero connection gives $h=0$, whereas the local connection
$A'=a_0dx$ gives
\begin{equation}
h^{A'}=a_0^2dx^2.
\label{eq:pb-connection-dependence-example}
\end{equation}
This example refutes a connection-independent interpretation of
$h_s^\omega$; it does not refute the gauge invariance of a fixed
connection--section pair.  \status{ESTABLISHED}

\section{Two channel geometries and their types}
\label{sec:pb-two-channels}

Restoring the channel labels gives the ordered pair
\begin{equation}
(h_{i,b}^{\omega_b},h_{i,m}^{\omega_m})
\in\Gamma(\operatorname{Sym}^2T^*\mathcal U)
\times\Gamma(\operatorname{Sym}^2T^*\mathcal U),
\label{eq:pb-two-channel-pair}
\end{equation}
and an analogous pair of cubic tensors.  This pair is determined by the two
selected statistical models, sections, and connections.  A single scalar
geometry is not part of those data.  \status{ESTABLISHED}

\hypothesisheading{Weighted product geometry}{hyp:pb-weighted-product-geometry}
If one joint geometry is required, choose constants $w_b,w_m>0$ and define
\begin{equation}
h^{\mathrm{prod}}
=w_bh_{i,b}^{\omega_b}+w_mh_{i,m}^{\omega_m}.
\label{eq:pb-weighted-channel-pullback}
\end{equation}
The weights fix the relative normalization of the two statistical channels.
Unit weights are a specialization rather than a consequence of the common
principal bundle.  \status{HYPOTHESIS}

Whenever weighted products at two different scales are compared, two further
declarations are consumed and neither follows from any fiberwise contraction
theorem.  Writing $\mu$ and $\bar\mu$ for the finite positive base measures
used to integrate the two tensors and $w_x,\bar w_x$ for the two sets of
channel weights, the comparison hypotheses are
\begin{equation}
\textbf{(X1)}\quad f_\#\mu=\bar\mu,
\qquad\qquad
\textbf{(X2)}\quad \bar w_x\circ f\leq w_x
\quad\text{for }x\in\{b,m\}.
\label{eq:pb-cross-scale-declarations}
\end{equation}
Under a channelwise comparison the two-channel difference is
\begin{equation}
h^{\mathrm{prod}}-f^*\bar h^{\mathrm{prod}}
=\sum_{x\in\{b,m\}}
\left[
w_x\left(h_x-f^*\bar h_x\right)
+\left(w_x-\bar w_x\circ f\right)f^*\bar h_x
\right],
\label{eq:pb-two-channel-comparison}
\end{equation}
so independently declared coarse weights with $(\bar w_x\circ f)>w_x$ break
positivity even when every channel has zero horizontal defect and a genuine
Markov fiber map.  Failure of (X1) reverses the integrated comparison just as directly: take
$\mathcal C=\bar{\mathcal C}$ the disjoint union of two copies of $\R$ with
$f=\operatorname{id}$, let $h$ vanish on the first copy and equal $dx^2$ on the
second, and let $\mu$ and $\bar\mu$ be unit point masses on the first and the
second copy respectively; the fine integrated tensor is zero and the coarse one
is one, with vanishing fiberwise defect throughout.  \status{ESTABLISHED}

\propositionheading{Radical of the weighted product}{prop:pb-product-radical}
Under \Cref{hyp:pb-regular-models,hyp:pb-weighted-product-geometry},
\begin{equation}
\operatorname{rad}h^{\mathrm{prod}}
=\ker D^{\omega_b}q_i\cap\ker D^{\omega_m}s_i.
\label{eq:pb-product-radical}
\end{equation}
Hence the two degenerate channel tensors can jointly be nondegenerate exactly
when their null distributions have zero intersection.  \status{ESTABLISHED}

\paragraph{Proof.}
For every $X$,
\[
h^{\mathrm{prod}}(X,X)
=w_bg_b^F(D^{\omega_b}q_iX,D^{\omega_b}q_iX)
+w_mg_m^F(D^{\omega_m}s_iX,D^{\omega_m}s_iX).
\]
Both summands are nonnegative and both weights are positive.  Their sum
vanishes exactly when both vertical jets vanish.  Polarization then identifies
the radical.  $\square$

The common principal group acts separately through $\widehat\rho_b$ and
$\widehat\rho_m$.  It supplies no bilinear cross tensor on
$V\mathcal E_b\times V\mathcal E_m$.  Such a cross term requires a separately
declared invariant bilinear form, and a comparison of channel points requires
$\Phi$ or $\widetilde\Phi$ from \Cref{sec:geo-cross-morphisms}.  Neither the
relative principal frame nor equal channel dimensions supplies either object.
\status{ESTABLISHED}

\section{Rank, radicals, and quotient geometry}
\label{sec:pb-rank-quotient}

\theoremheading{Constant-rank quotient theorem}{thm:pb-pullback-rank-quotient}
Let $s:\mathcal U\to E$ satisfy
\Cref{hyp:pb-regular-models}.  Pointwise,
\begin{equation}
\operatorname{rad}h_s^\omega
=\ker D^\omega s,
\qquad
\operatorname{rank}h_s^\omega
=\operatorname{rank}D^\omega s.
\label{eq:pb-pullback-radical-rank}
\end{equation}
If $D^\omega s$ has constant rank on $\mathcal U$, then
$K_s^\omega=\ker D^\omega s$ is a smooth subbundle and
\begin{equation}
Q_s^\omega=T\mathcal U/K_s^\omega,
\label{eq:pb-pullback-quotient-bundle}
\end{equation}
has a positive-definite metric and a symmetric cubic tensor defined by
\begin{align}
\bar h_s^\omega([X],[Y])
&=h_s^\omega(X,Y),
\label{eq:pb-pullback-quotient-metric}\\
\bar c_s^\omega([X],[Y],[Z])
&=c_s^\omega(X,Y,Z).
\label{eq:pb-pullback-quotient-amari}
\end{align}
The induced map $Q_s^\omega\to\operatorname{im}D^\omega s$ is an isometric
vector-bundle isomorphism.  \status{ESTABLISHED}

\paragraph{Proof.}
If $D^\omega sX=0$, then \eqref{eq:pb-fisher-pullback} gives
$h_s^\omega(X,Y)=0$ for all $Y$.  Conversely, if $X$ lies in the radical,
then
\[
0=h_s^\omega(X,X)
=g^F(D^\omega sX,D^\omega sX),
\]
and positive definiteness of $g^F$ forces $D^\omega sX=0$.  This proves
\eqref{eq:pb-pullback-radical-rank}.  The applicable constant-rank statement is
the one for morphisms of smooth vector bundles over a fixed base, not the
constant-rank theorem for smooth maps of manifolds: $D^\omega s$ is a
vector-bundle morphism $T\mathcal U\to s^*VE$ over $\operatorname{id}$, and a
vector-bundle morphism of locally constant rank has smooth kernel and image
subbundles by local adaptation of frames.  Replacing an argument of
\eqref{eq:pb-pullback-quotient-metric} or
\eqref{eq:pb-pullback-quotient-amari} by a representative differing by an
element of $K_s^\omega$ leaves its vertical jet unchanged.  The quotient
tensors are therefore well defined.  The factored map induced by
$D^\omega s$ is injective onto its image and preserves the displayed tensors;
the quotient metric is positive definite.  $\square$

The quotient in \Cref{thm:pb-pullback-rank-quotient} is a vector bundle.  It
is not automatically the tangent bundle of a quotient manifold.  For that
interpretation, require a smooth Hausdorff quotient manifold
$\overline{\mathcal U}$ and a surjective submersion
$q_K:\mathcal U\to\overline{\mathcal U}$ whose fibers are exactly the null
leaves and whose tangent kernel is
$\ker Tq_K=K_s^\omega$.  Consequently, the null distribution must be
involutive,
\begin{equation}
D^\omega s([X,Y])=0
\quad
\text{for all }X,Y\in\Gamma(K_s^\omega),
\label{eq:pb-null-involutivity}
\end{equation}
but involutivity alone supplies neither the quotient manifold nor the
surjective-submersion quotient map. \status{ESTABLISHED}

Separately, the tensors must be basic for the null foliation.  Their contraction with a null
vector already vanishes by construction, but invariance along null leaves is
an additional requirement:
\begin{equation}
\mathcal L_Zh_s^\omega=0,
\qquad
\mathcal L_Zc_s^\omega=0
\quad\text{for every }Z\in\Gamma(K_s^\omega).
\label{eq:pb-null-basicness}
\end{equation}
Without this condition, even a smooth Hausdorff leaf-space quotient need not
carry tensors whose pullbacks are $h_s^\omega$ and $c_s^\omega$.
\status{ESTABLISHED}

Basicness of the quadratic tensor is checkable directly on the jet, but the
bundle-valued one-form $D^\omega s$ has no ordinary Lie derivative.  Choose a
$g^F$-compatible auxiliary connection $\nabla$ on $s^*VE$ and, for a
$s^*VE$-valued one-form $\alpha$, write
\begin{equation}
(\mathcal L_Z^\nabla\alpha)(X)
=\nabla_Z(\alpha(X))-\alpha([Z,X]).
\label{eq:pb-covariant-lie-derivative}
\end{equation}
Then for $Z\in\Gamma(K_s^\omega)$ and arbitrary $X,Y$, the Leibniz rule for
the ordinary Lie derivative of $h_s^\omega$ gives
\begin{equation}
(\mathcal L_Zh_s^\omega)(X,Y)
=g^F\left((\mathcal L_Z^\nabla D^\omega s)(X),D^\omega sY\right)
+g^F\left(D^\omega sX,(\mathcal L_Z^\nabla D^\omega s)(Y)\right).
\label{eq:pb-basicness-criterion}
\end{equation}
Metric compatibility supplies this identity directly, while
$\iota_Zh_s^\omega=0$ is already automatic from $Z\in\ker D^\omega s$.
If $\nabla'=\nabla+B$ is another $g^F$-compatible connection, then $B_Z$ is
$g^F$-skew-adjoint, so the changes in the two displayed covariant summands
cancel.  Thus the displayed equality with the ordinary, canonical
$\mathcal L_Zh_s^\omega$ is independent of the selected auxiliary metric
connection, even though neither covariant summand is separately canonical.
\status{ESTABLISHED}

\propositionheading{A constant-rank nonintegrable radical}{prop:pb-contact-null-counterexample}
There is a regular one-dimensional statistical fiber, a smooth connection,
and a section over $\R^3$ for which $h_s^\omega$ has constant rank one while
$\ker D^\omega s$ is not involutive.  \status{ESTABLISHED}

\paragraph{Proof.}
Use the normal location family and the trivial principal translation bundle
over $\R^3$ with coordinates $(x,y,z)$.  In the constant section $\mu=0$,
choose the local connection form
\begin{equation}
\alpha=\mathrm{d}z-x\mathrm{d}y.
\label{eq:pb-contact-connection}
\end{equation}
Then $D^\omega s=\alpha$ and
\begin{equation}
h_s^\omega=(\mathrm{d}z-x\mathrm{d}y)^2.
\label{eq:pb-contact-pullback}
\end{equation}
The rank is one everywhere, but
\begin{equation}
\alpha\wedge\mathrm{d}\alpha
=-\mathrm{d}z\wedge\mathrm{d}x\wedge\mathrm{d}y\ne0.
\label{eq:pb-contact-frobenius}
\end{equation}
The Frobenius criterion therefore shows that $\ker\alpha$ is a contact
distribution rather than a foliation.  $\square$

Constant rank is also necessary for the smooth quotient-bundle conclusion.
For the same normal location family on $\mathcal C=\R$, with the zero
connection and $s(x)=\mathcal N(x^2,1)$,
\begin{equation}
h_s=4x^2dx^2,
\label{eq:pb-rank-jump-example}
\end{equation}
has rank one away from the origin and rank zero at the origin.  The pointwise
quotients do not assemble into a vector bundle across that point.
\status{ESTABLISHED}

\section{Divergence jets and transported local comparison}
\label{sec:pb-divergence-jets}

Let $\mathscr D:\mathcal B\times\mathcal B\to\R$ be a $C^3$ contrast in a
neighborhood of the diagonal: $\mathscr D(p,q)\geq0$ and
$\mathscr D(p,q)=0$ exactly on the diagonal in that neighborhood.  For tangent
vectors below, assume additionally that $\mathscr D$ is invariant under the
simultaneous represented action,
$\mathscr D(\widehat\rho(g)p,\widehat\rho(g)q)=\mathscr D(p,q)$, so that it
defines an unambiguous comparison of two points in one associated-bundle
fiber.  Without this invariance, the following jets are only local-frame
objects.  Assume also that the mixed second form defined below is positive
definite on the selected regular tier.  A contrast such as
$\mathscr D(\mu,\nu)=(\mu-\nu)^4$ is positive off the diagonal but has a
zero diagonal Hessian; it yields degenerate jets rather than a divergence-
induced dualistic geometry. \status{HYPOTHESIS}

For tangent
vectors $u,v,w\in T_p\mathcal B$, let $(u,0)$ and $(0,u)$ denote their lifts
to the first and second factors.  In one common local chart, extend these
vectors as constant coordinate vector fields and define the mixed diagonal
jets
\begin{align}
g_{\mathscr D}(u,v)
&=-[(u,0)(0,v)\mathscr D]_{(p,p)},
\label{eq:pb-divergence-metric-jet}\\
\Gamma_{\mathscr D}(u,v,w)
&=-[(u,0)(v,0)(0,w)\mathscr D]_{(p,p)},
\label{eq:pb-divergence-first-connection-jet}\\
\Gamma_{\mathscr D}^*(u,v,w)
&=-[(0,u)(0,v)(w,0)\mathscr D]_{(p,p)},
\label{eq:pb-divergence-dual-connection-jet}\\
\mathcal T_{\mathscr D}
&=\Gamma_{\mathscr D}-\Gamma_{\mathscr D}^*.
\label{eq:pb-divergence-amari-jet}
\end{align}
The individual third-order arrays are connection coefficients in a chart;
their difference is a cubic tensor.  The sign convention in
\eqref{eq:pb-divergence-amari-jet} is chosen so that forward KL on an
exponential family gives the score-third-moment tensor
\eqref{eq:pb-fiber-amari}.  \status{DEFINITION}

\propositionheading{KL divergence jets}{prop:pb-kl-divergence-jets}
On a regular statistical model satisfying the integrability hypotheses of
\Cref{hyp:pb-regular-models}, the mixed second jet of
$\KL(p\Vert q)$ is the Fisher metric, and the difference of the two mixed
third jets is the Amari--Chentsov tensor.  \status{ESTABLISHED}

\paragraph{Proof.}
Use a local parameter $\theta$ for the first argument and $\eta$ for the
second.  Write $\ell_i=\partial_i\log p_\theta$ and let primed derivatives
act on $\eta$.  Normalization gives $\E_\theta[\ell_i]=0$, while
$\partial_i p_\theta=p_\theta\ell_i$ and
$\partial_i\partial_jp_\theta
=p_\theta(\ell_i\ell_j+\partial_i\ell_j)$.  Direct differentiation at
$\eta=\theta$ therefore gives
\begin{align}
-\partial_i\partial_j'\KL(p_\theta\Vert p_\eta)
&=\E_\theta[\ell_i\ell_j],\\
-\partial_i\partial_j\partial_k'\KL(p_\theta\Vert p_\eta)
&=\E_\theta[(\ell_i\ell_j+\partial_i\ell_j)\ell_k],\\
-\partial_i'\partial_j'\partial_k\KL(p_\theta\Vert p_\eta)
&=\E_\theta[\ell_k\partial_i\ell_j].
\label{eq:pb-kl-jet-calculation}
\end{align}
The score Hessian is symmetric in the chosen smooth chart, so subtracting
the last line from the second cancels its score-derivative term and leaves
$\E_\theta[\ell_i\ell_j\ell_k]$.  Polarization gives the asserted formulas
for arbitrary $u,v,w$.  Differentiation under the integral is licensed by
\Cref{hyp:pb-regular-models}; these are the standard divergence-induced
dualistic structures of information geometry
\citep{AmariNagaoka2000,AyJostLeSchwachhofer2017}.  $\square$

For a smooth base curve $\gamma$ with $\gamma(0)=c$, let
$\Omega^\omega_{0\to\epsilon}$ be the channel transport along
$\gamma|_{[0,\epsilon]}$ and bring the section value back to the initial fiber by
\begin{equation}
\widehat s_\gamma(\epsilon)
=(\Omega^\omega_{0\to\epsilon})^{-1}s(\gamma(\epsilon)).
\label{eq:pb-transported-section-curve}
\end{equation}
Then
\begin{equation}
\left.\frac{d}{d\epsilon}\widehat s_\gamma(\epsilon)\right|_{\epsilon=0}
=D^\omega s(\dot\gamma(0)).
\label{eq:pb-transported-section-velocity}
\end{equation}
\status{ESTABLISHED}

\corollaryheading{Quadratic transported divergence}{cor:pb-transported-divergence-quadratic}
Let $g_{\mathscr D}$ be the metric induced by
\eqref{eq:pb-divergence-metric-jet} and set
$h_{s,\mathscr D}^\omega=(D^\omega s)^*g_{\mathscr D}$.  Then
\begin{equation}
\mathscr D(s(c),\widehat s_\gamma(\epsilon))
=\frac{\epsilon^2}{2}h_{s,\mathscr D}^\omega
(\dot\gamma(0),\dot\gamma(0))
+O(\epsilon^3).
\label{eq:pb-transported-divergence-expansion}
\end{equation}
If $g_{\mathscr D}=g^F$, in particular for the regular KL contrast of
\Cref{prop:pb-kl-divergence-jets}, then
$h_{s,\mathscr D}^\omega=h_s^\omega$.
\status{ESTABLISHED}

\paragraph{Proof.}
A contrast and its first derivative vanish on the diagonal.  Its quadratic
Taylor term in the second argument is the induced metric.  Substituting
\eqref{eq:pb-transported-section-velocity} gives
\eqref{eq:pb-transported-divergence-expansion}.  $\square$

The coefficient of $\epsilon^3$ in a one-sided expansion such as
\eqref{eq:pb-transported-divergence-expansion} is not determined by
$c_{s,\mathscr D}^\omega:=(D^\omega s)^*\mathcal T_{\mathscr D}$ alone.
It also contains the covariant acceleration of
$\widehat s_\gamma$ and the connection coefficient selected by the one-sided
divergence jet.  The divergence-induced cubic is recovered from the
difference of the two mixed third jets in
\eqref{eq:pb-divergence-amari-jet}, not from an arbitrary one-sided cubic
Taylor coefficient.  For regular KL, it is the Amari--Chentsov tensor and
$c_{s,\mathscr D}^\omega=c_s^\omega$. \status{ESTABLISHED}

\section{Vertical information length}
\label{sec:pb-curve-taxonomy}

For a smooth curve $\Gamma:J\to E$ over $\gamma=\varpi\circ\Gamma$, define its
connection-relative vertical velocity by
\begin{equation}
v^\omega(\Gamma)
=\operatorname{ver}^{\omega}(\dot\Gamma)
=\dot\Gamma-H^\omega_\Gamma\dot\gamma,
\label{eq:pb-curve-vertical-velocity}
\end{equation}
and its vertical Fisher length by
\begin{equation}
L^\omega(\Gamma)
=\int_J\sqrt{g^F(v^\omega(\Gamma),v^\omega(\Gamma))}d\lambda.
\label{eq:pb-vertical-fisher-length}
\end{equation}
Here $\lambda$ is an arbitrary curve parameter rather than a declared
physical time.  \status{DEFINITION}

If $\omega'=\omega+a$, define along the curve
\begin{equation}
R_a^\Gamma(\dot\gamma)
=\vartheta_{\Gamma(\lambda)}
\left(a_{\gamma(\lambda)}(\dot\gamma(\lambda))\right).
\label{eq:pb-curve-connection-change}
\end{equation}
With the convention of \eqref{eq:pb-jet-connection-change}, the exact curve
identity is
$v^{\omega'}(\Gamma)=v^\omega(\Gamma)+R_a^\Gamma(\dot\gamma)$.
\status{ESTABLISHED}

\propositionheading{Vertical, horizontal, and section curves}{prop:pb-curve-taxonomy}
The length $L^\omega$ is invariant under orientation-preserving $C^1$
reparameterization with positive derivative.  If $\Gamma$ is vertical, then
$L^\omega(\Gamma)$ is independent of the connection.  If $\Gamma$ is
$\omega$-horizontal, then
$L^\omega(\Gamma)=0$.  For a section curve $\Gamma=s\circ\gamma$,
\begin{equation}
L^\omega(s\circ\gamma)
=\int_J\sqrt{h_s^\omega(\dot\gamma,\dot\gamma)}d\lambda.
\label{eq:pb-section-curve-length}
\end{equation}
For a mixed curve, changing $\omega$ to $\omega'$ changes the vertical
velocity by $R_a^\Gamma(\dot\gamma)$ in the convention of
\eqref{eq:pb-jet-connection-change}; its length is therefore generally
connection dependent.  \status{ESTABLISHED}

\paragraph{Proof.}
Under a reparameterization $\lambda=\phi(r)$, vertical velocity is multiplied
by $\phi'(r)$, and positive homogeneity of the square root converts that
factor into the Jacobian in the one-dimensional change-of-variables formula.
For a vertical curve, $\dot\gamma=0$, so the horizontal correction vanishes
for every connection.  Horizontality is the equation
$v^\omega(\Gamma)=0$.  If $\Gamma=s\circ\gamma$, then
$v^\omega(\Gamma)=D^\omega s(\dot\gamma)$, and
\eqref{eq:pb-section-curve-length} follows from
\eqref{eq:pb-fisher-pullback}.  The connection-change statement is
\eqref{eq:pb-jet-connection-change} along the curve.  $\square$

The functional \eqref{eq:pb-vertical-fisher-length} measures only vertical
statistical change.  It assigns zero length to horizontal displacement and
does not define a metric on $TE$ without a separately chosen base metric.  It
also selects neither an inference orbit nor an orientation.  Those require a
separately declared configuration manifold together with an objective
differential and a mobility or metric sharp map that defines its vector field.
\status{NOT-CLAIMED}

A base curve $\gamma$ orders contexts, and $s\circ\gamma$ evaluates one fixed
section along that ordering; neither object is an agent inference history.
Updating at one fixed $c$ is instead a vertical curve in the relevant
statistical fiber, while evolution of a complete contextual agent is a curve
in a separately declared section configuration space.
\status{DEFINITION}

The finite-dimensional length above applies pointwise or on a declared
finite-design product.  It does not construct a continuum section-space
length or a probability law on a continuum space of sections, because no base
measure, channel weighting, or gauge-quotient rule has been declared here.
\status{NOT-CLAIMED}

\section{Coarse-morphism naturality}
\label{sec:pb-coarse-naturality}

Consider smooth associated bundles $E\to\mathcal C$ and
$\bar E\to\bar{\mathcal C}$ with connections $\omega$ and $\bar\omega$.
Let $f:\mathcal C\to\bar{\mathcal C}$ and let
$\Psi:E\to\bar E$ be a smooth bundle morphism over $f$.  Suppose sections
$s$ and $\bar s$ are related as follows.  Before that relation is imposed,
$\Psi\circ s$ is only a section of the pullback bundle $f^*\bar E$ along
$f$; a fine section does not by itself define a section on
$\bar{\mathcal C}$.  \status{DEFINITION}

We adopt the descent hypothesis that there exists a smooth coarse section
$\bar s$ satisfying
\begin{equation}
\Psi\circ s=\bar s\circ f.
\label{eq:pb-coarse-related-sections}
\end{equation}
For a surjective submersion $f$, this excludes fine sections whose
$\Psi$-images vary along a fiber of $f$. \status{HYPOTHESIS}

Write $T^V\Psi:VE\to V\bar E$ for the vertical differential. Following the
cross-channel defect of \eqref{eq:geo-nonlinear-defect-phi}, define
\begin{equation}
(\mathcal D\Psi)_e(X)
=A_\Psi(e;X)
:=\operatorname{ver}^{\bar\omega}
\left(T_e\Psi(H_e^\omega X)\right)
=T_e\Psi(H_e^\omega X)
-H_{\Psi(e)}^{\bar\omega}(T_cfX).
\label{eq:pb-coarse-horizontal-defect}
\end{equation}
The sign convention is fixed and load bearing: the transported fine horizontal
lift comes first and the coarse horizontal lift is subtracted.  The difference
is vertical because
$T\bar\varpi\circ T\Psi(H^\omega_eX)=T_cfX=T\bar\varpi\circ H^{\bar\omega}(T_cfX)$,
and the second displayed form follows because the $\bar\omega$-horizontal part
of $T\Psi(H^\omega_eX)$ is exactly $H^{\bar\omega}_{\Psi(e)}(T_cfX)$.  Its type
is
$A_\Psi\in\Gamma(E;\varpi^*T^*\mathcal C\otimes\Psi^*V\bar E)$, and along a
section
$A_\Psi(s;\cdot)\in\Gamma(\mathcal U;T^*\mathcal U\otimes f^*\bar s^*V\bar E)$.
The abbreviation $A_\Psi$ is used below wherever the vertical values must be
paired with $\bar g^F$.  \status{DEFINITION}

\theoremheading{Sharp section descent}{thm:pb-section-descent}
Let $f$ be a surjective smooth submersion, let $\Psi$ be a smooth bundle
morphism over $f$, and let $Q\in\Gamma(\mathcal C,E)$.  Consider the three
conditions: (P1) there exists $\bar Q\in\Gamma(\bar{\mathcal C},\bar E)$ with
$\Psi\circ Q=\bar Q\circ f$; (P2) $\Psi\circ Q$ is constant on each fiber of
$f$; and
\begin{equation}
\textbf{(P3)}\qquad
T^V\Psi\left(D^\omega Q(X)\right)+A_\Psi\left(Q(c);X\right)=0
\quad\text{for every }c\text{ and every }X\in\ker T_cf.
\label{eq:pb-projectability-criterion}
\end{equation}
Then (P1) and (P2) are equivalent, both imply (P3), and (P3) implies (P2) when
the fibers of $f$ are connected.  Under (P2) the descended $\bar Q$ is unique,
is automatically smooth, and is automatically a section.  \status{ESTABLISHED}

\paragraph{Proof.}
That (P1) implies (P2) is immediate.  Conversely, a surjective smooth
submersion is a smooth quotient map, so a smooth map constant on its fibers
descends uniquely through it and the descended factor is smooth; this is where
smoothness of $\bar Q$ becomes a theorem rather than an obligation.
Surjectivity gives uniqueness, and
$\bar\varpi\bar Qf=\bar\varpi\Psi Q=f\varpi Q=f$ with $f$ surjective gives
$\bar\varpi\bar Q=\operatorname{id}$.  For (P2) implying (P3), take
$X\in\ker T_cf$, split $T_cQX=H^\omega_{Q(c)}X+D^\omega Q(X)$, apply $T\Psi$,
and use $T_cfX=0$ to annihilate the coarse horizontal lift; what remains is
$T(\Psi Q)(X)=A_\Psi(Q(c);X)+T^V\Psi(D^\omega Q(X))$, which fiber constancy
forces to vanish.  The same computation shows that (P3) makes $T(\Psi Q)$
annihilate $\ker Tf$, and a smooth map with vanishing derivative along a
connected embedded submanifold is constant on it.  $\square$

The submersion hypothesis is load bearing rather than decorative.  With
$\mathcal C=\bar{\mathcal C}=\R$, $f(x)=x^3$ a smooth bijection that fails to be
a submersion at the origin, the trivial bundle with fiber
$\{\mathcal N(\mu,1)\}$, $\Psi$ the identity on fibers, and
$Q(x)=\mathcal N(x,1)$, the unique descent $\bar Q(y)=\mathcal N(y^{1/3},1)$ is
continuous but not differentiable at $y=0$.  \status{ESTABLISHED}

\propositionheading{A pointwise bundle morphism is not a map on configurations}{prop:pb-nonfunctorial-descent}
Suppose $f$ is a surjective submersion, $\ker T_{c_0}f\ne0$ at some $c_0$,
there are $e_0\in E_{c_0}$ and $w\in V_{e_0}E$ with
$T^V\Psi(w)\ne0$, and a global smooth section $Q_0$ satisfies
$Q_0(c_0)=e_0$.  Then the projectable set
$\Gamma_{\mathrm{proj}}(\Psi)$ of sections satisfying (P2) is a proper subset
of $\Gamma(\mathcal C,E)$.  \status{ESTABLISHED}

\paragraph{Proof.}
Choose a chart in which $\ker Tf$ is spanned by
$\partial_1,\dots,\partial_k$ with $k\geq1$, and use the declared global
section $Q_0$ through $e_0$.  If $Q_0$ already fails
(P2) there is nothing to prove.  Otherwise let $W$ be a smooth vertical field
in that local trivialization of $E$ with $W(Q_0(c_0))=w$ and local flow
$\Phi^W$, take
$\chi\in C^\infty_c$ supported inside this trivializing chart with
$\chi(c_0)=0$ and $\partial_1\chi(c_0)\ne0$, and define the global smooth
section $Q_\epsilon$ by
$Q_\epsilon(c)=\Phi^W_{\epsilon\chi(c)}(Q_0(c))$ on the support chart and
$Q_\epsilon=Q_0$ off its support.  Since
$\chi(c_0)=0$, both $A_\Psi(Q_\epsilon(c_0);\partial_1)$ and $T^V\Psi$ at that
point are independent of $\epsilon$, so differentiating the left-hand side of
\eqref{eq:pb-projectability-criterion} at $\epsilon=0$ gives
$\partial_1\chi(c_0)T^V\Psi(w)\ne0$ and (P3) fails for small nonzero
$\epsilon$.  $\square$

A concrete witness is the total collapse $\mathcal C=S^1$,
$\bar{\mathcal C}=\{\ast\}$, which is a surjective submersion, with trivial
bundles carrying the unit-variance normal location fiber, $\Psi$ the identity
on fibers, and $Q(x)=\mathcal N(\sin x,1)$.  Here $A_\Psi=0$ and
$\ker Tf=TS^1$, while
$T^V\Psi(D^\omega Q\partial_x)=\cos x\partial_\mu$ is not identically zero, so
no coarse section exists.  In the linear realization with configuration space
$L^2(S^1)$ the descendable set is exactly the constants: a closed subspace of
infinite codimension and empty interior.  This statement is scoped to a
collapsing $f$; if $f$ is a diffeomorphism every section descends, so no
unscoped claim that a bundle morphism induces a configuration map nowhere is
available.  \status{ESTABLISHED}

\theoremheading{Covariant first-jet chain rule}{thm:pb-covariant-jet-naturality}
Under \eqref{eq:pb-coarse-related-sections},
\begin{equation}
D^{\bar\omega}\bar s(TfX)
=T^V\Psi(D^\omega sX)+(\mathcal D\Psi)_{s(c)}(X).
\label{eq:pb-covariant-jet-chain-rule}
\end{equation}
Hence the square of covariant first jets commutes exactly when
$\mathcal D\Psi$ vanishes along $s$.  In that case, every vertical covariant
tensor $\bar\tau$ on $\bar E$ satisfies
\begin{equation}
f^*[(D^{\bar\omega}\bar s)^*\bar\tau]
=(D^\omega s)^*[(T^V\Psi)^*\bar\tau].
\label{eq:pb-covariant-tensor-naturality}
\end{equation}
\status{ESTABLISHED}

\paragraph{Proof.}
Differentiate \eqref{eq:pb-coarse-related-sections} and decompose
$TsX=H_s^\omega X+D^\omega sX$.  Applying $T\Psi$ to the horizontal term
gives the target horizontal lift plus
$(\mathcal D\Psi)_{s(c)}(X)$ by
\eqref{eq:pb-coarse-horizontal-defect}; applying it to the vertical term
gives $T^V\Psi(D^\omega sX)$.  The vertical part of the other side is
$D^{\bar\omega}\bar s(TfX)$, which proves
\eqref{eq:pb-covariant-jet-chain-rule}.  When the defect vanishes, applying
$\bar\tau$ to the intertwined vertical jets proves
\eqref{eq:pb-covariant-tensor-naturality}.  $\square$

The base comparison of the two Fisher pullbacks can now be written exactly,
with no compatibility hypothesis of any kind.  Along $s$ write
$u_X=D^\omega sX$ for the fine vertical jet, $L=T^V\Psi$ for the vertical
differential, $a_X=A_\Psi(s(c);X)$ for the horizontal defect, and
$\bar u_X=D^{\bar\omega}\bar s(T_cfX)=Lu_X+a_X$ for the coarse vertical jet.
With the vertical Fisher defect $\Delta_F^\Psi$ of
\eqref{eq:pb-vertical-fisher-defect} below, set
\begin{align}
\delta_\Psi(X,Y)&=\Delta_F^\Psi(u_X,u_Y),
\label{eq:pb-base-fisher-defect}\\
\mathcal X_\Psi(X,Y)&=\bar g^F(Lu_X,a_Y)+\bar g^F(a_X,Lu_Y),
\label{eq:pb-cross-anomaly-tensor}\\
\mathcal Q_\Psi(X,Y)&=\bar g^F(a_X,a_Y).
\label{eq:pb-quadratic-anomaly-tensor}
\end{align}

\theoremheading{Exact signed base Fisher comparison}{thm:pb-signed-base-comparison}
Under \eqref{eq:pb-coarse-related-sections} alone,
\begin{equation}
h_s^\omega-f^*\bar h_{\bar s}^{\bar\omega}
=\delta_\Psi-\mathcal X_\Psi-\mathcal Q_\Psi
\label{eq:pb-signed-base-comparison}
\end{equation}
as symmetric two-tensors on the fine domain.  Here $\mathcal X_\Psi$ is
symmetric and sign indefinite while $\mathcal Q_\Psi\succeq0$, so both retained
tensors enter with a minus sign and neither may be summarized as one vertical
mismatch term.  \status{ESTABLISHED}

\paragraph{Proof.}
By \eqref{eq:pb-covariant-jet-chain-rule},
$f^*\bar h(X,Y)=\bar g^F(Lu_X+a_X,Lu_Y+a_Y)
=\bar g^F(Lu_X,Lu_Y)+\mathcal X_\Psi(X,Y)+\mathcal Q_\Psi(X,Y)$, and by
\eqref{eq:pb-vertical-fisher-defect},
$\bar g^F(Lu_X,Lu_Y)=g^F(u_X,u_Y)-\Delta_F^\Psi(u_X,u_Y)$.  Subtracting from
$h_s^\omega(X,Y)=g^F(u_X,u_Y)$ gives
\eqref{eq:pb-signed-base-comparison}.  $\square$

\theoremheading{Exact signed positivity criterion}{thm:pb-signed-positivity-criterion}
Assume in addition the Markov hypotheses of
\Cref{thm:pb-pullback-fisher-defect}, so that $\delta_\Psi\succeq0$.  Then at a
fixed $c$ the following are equivalent: the base comparison
$h_s^\omega-f^*\bar h_{\bar s}^{\bar\omega}$ is positive semidefinite; the
coarse jet is no longer than the fine jet, meaning
$\lVert D^{\bar\omega}\bar s(T_cfX)\rVert_{\bar g^F}
\leq\lVert D^\omega sX\rVert_{g^F}$ for every $X$; and
\begin{equation}
2\bar g^F\left(T^V\Psi D^\omega sX,A_\Psi(s(c);X)\right)
+\left\lVert A_\Psi(s(c);X)\right\rVert^2_{\bar g^F}
\leq\delta_\Psi(X,X)
\quad\text{for every }X.
\label{eq:pb-positivity-criterion}
\end{equation}
Consequently vanishing of the horizontal defect is \emph{sufficient} for base
positivity and is \emph{not necessary}: whenever the cross term
$\bar g^F(Lu_X,a_X)$ is negative and large enough in modulus,
\eqref{eq:pb-positivity-criterion} holds strictly with $a_X\ne0$.  The
triangle-inequality margin
$\lVert a_X\rVert_{\bar g^F}\leq\sqrt{h(X,X)}-\sqrt{h(X,X)-\delta_\Psi(X,X)}$
is likewise sufficient and not necessary.  \status{ESTABLISHED}

\paragraph{Proof.}
The first two conditions are \eqref{eq:pb-signed-base-comparison} read on the
diagonal, using $h(X,X)=\lVert u_X\rVert^2_{g^F}$ and
$f^*\bar h(X,X)=\lVert Lu_X+a_X\rVert^2_{\bar g^F}$.  Expanding that square and
substituting
$\lVert Lu_X\rVert^2_{\bar g^F}=\lVert u_X\rVert^2_{g^F}-\delta_\Psi(X,X)$
gives \eqref{eq:pb-positivity-criterion}.  $\square$

Both failure directions are realized on the normal location fiber, and the
arithmetic is exact.  Take $\mathcal C=\bar{\mathcal C}=\R$, $f$ the identity,
$\mathcal B=\{\mathcal N(\mu,1)\}$ with $g^F=1$, the sample kernel
$N(x,\cdot)=\mathcal N(x,1)$ so that
$\bar{\mathcal B}=\{\mathcal N(\mu,2)\}$ with $\bar g^F=\tfrac12$ and
$\Delta_F^\Psi=\tfrac12$, the section $s(x)=\mathcal N(x,1)$, zero source
connection, and a target connection contributing $a_X=b\partial_\mu$.  Then
$h_s^\omega-f^*\bar h_{\bar s}^{\bar\omega}=1-\tfrac12(1+b)^2$, which equals
$\tfrac{79}{200}>0$ at $b=\tfrac1{10}$ and $\tfrac{23}{25}>0$ at $b=-\tfrac35$,
both with nonzero horizontal defect, and equals $-\tfrac18<0$ at
$b=\tfrac12$.  At $b=-\tfrac35$ the margin above is violated, since
$\lVert a_X\rVert_{\bar g^F}=\tfrac{3}{5\sqrt2}$ exceeds
$1-\tfrac1{\sqrt2}$, so the margin is genuinely only sufficient.  Strict
negativity survives even at zero information loss: with the identity kernel, so
$\Delta_F^\Psi=0$, constant related sections, source horizontal field
$\partial_x$, and target horizontal field $\partial_x+a\partial_\mu$ with
$a\ne0$, one gets $u=0$, $a_X=-a\partial_\mu$, and
$h_s^\omega-f^*\bar h_{\bar s}^{\bar\omega}=-a^2dx^2\prec0$.
\status{ESTABLISHED}

The horizontal defect of an induced morphism is not an arbitrary vertical
field.  Suppose $\Psi$ is induced by an equivariant principal scale map
$\mathcal P$ over $f$ with Lie-group homomorphism $\kappa$ and law-fiber map
$q$, in the sense of \eqref{eq:rg-associated-scale-map}.  Define the
\emph{scale-connection defect form}
\begin{equation}
\mathfrak A_{\mathcal P}
=\mathcal P^*\bar\omega-\mathrm d\kappa\circ\omega
\in\Omega^1(P,\bar{\mathfrak g}).
\label{eq:pb-scale-connection-defect}
\end{equation}

\theoremheading{Isotropy criterion for a vanishing horizontal defect}{thm:pb-isotropy-criterion}
The form $\mathfrak A_{\mathcal P}$ is horizontal and
$\operatorname{Ad}\circ\kappa$-equivariant, hence descends to an element of
$\Omega^1(\mathcal C;f^*\operatorname{Ad}\bar P)$, and
\begin{equation}
A_\Psi(e;X)
=\vartheta_{\Psi(e)}\left(\mathfrak A_{\mathcal P}(X)\right),
\label{eq:pb-anomaly-fundamental-field}
\end{equation}
so the horizontal defect is a fundamental vertical field determined by a
$\bar{\mathfrak g}$-valued one-form on the base rather than by the fiber point.
Consequently, for related sections satisfying
$\Psi\circ s=\bar s\circ f$, the horizontal defect vanishes along the fine
section, meaning
$A_\Psi(s(c);X)=0$ for every $c\in\mathcal C$ and
$X\in T_c\mathcal C$, if and only if
\begin{equation}
\mathfrak A_{\mathcal P}(X)\in\bar{\mathfrak g}_{\bar s(f(c))}
\quad\text{for every }c\text{ and every }X,
\label{eq:pb-isotropy-criterion}
\end{equation}
the isotropy subalgebra of the coarse section value under
$\widehat{\bar\rho}$.  The principal-level identity $\mathfrak A_{\mathcal P}=0$
is sufficient, and is necessary only when the represented action is
infinitesimally effective at $\bar s(f(c))$.  Moreover
$\mathcal Q_\Psi=\mathfrak a^*\mathfrak k_{\bar s(f(c))}$ with
$\mathfrak k_{\bar\beta}(\xi,\eta)
=\bar g^F(\bar\zeta_\xi\bar\beta,\bar\zeta_\eta\bar\beta)\succeq0$, whose
radical is exactly $\bar{\mathfrak g}_{\bar\beta}$.  \status{ESTABLISHED}

\paragraph{Proof.}
Evaluating $\mathcal P^*\bar\omega$ on the fundamental field of $\xi$ gives
$\bar\omega(\bar\zeta_{\mathrm d\kappa\xi})=\mathrm d\kappa\xi$, so the
difference annihilates vertical vectors; equivariance follows from
$R_g^*\mathcal P^*\bar\omega
=\operatorname{Ad}_{\kappa(g)^{-1}}\mathcal P^*\bar\omega$ together with
$\mathrm d\kappa\circ\operatorname{Ad}_{g^{-1}}
=\operatorname{Ad}_{\kappa(g)^{-1}}\circ\mathrm d\kappa$.  In local
trivializations write $\Psi(c,\beta)=(f(c),\psi_c\beta)$ with
$\psi_c=\widehat{\bar\rho}(\varsigma(c))\circ q$.  Using
$H^\omega_{(c,\beta)}X=(X,-\zeta_{A(X)}\beta)$ and the differentiated
intertwining $Tq\circ\zeta_\xi=\bar\zeta_{\mathrm d\kappa(\xi)}\circ q$,
obtained from \eqref{eq:rg-scale-intertwiner} at $g=\exp(t\xi)$, gives
$A_\Psi(X)=\bar\zeta_{\mathfrak a(X)}(\psi_c\beta)$ with
$\mathfrak a=\operatorname{Ad}_\varsigma(u^*\mathfrak A_{\mathcal P})$, which is
\eqref{eq:pb-anomaly-fundamental-field}.  A fundamental field vanishes at
$\bar\beta$ exactly when its generator lies in $\bar{\mathfrak g}_{\bar\beta}$,
and conjugation by $\operatorname{Ad}_\varsigma$ moves both the generator and
the isotropy algebra together, which gives
\eqref{eq:pb-isotropy-criterion}.  The last statement is
\eqref{eq:pb-anomaly-fundamental-field} substituted into
\eqref{eq:pb-quadratic-anomaly-tensor}.  $\square$

\Cref{eq:pb-isotropy-criterion} is the definition of connection compatibility
used throughout this manuscript.  No weaker or informal reading is intended,
and no result below relies on a compatibility notion that has not been reduced
to it.  \status{DEFINITION}

\theoremheading{Ordered composition of horizontal defects}{thm:pb-anomaly-composition}
For composable morphisms
$E_0\xrightarrow{\Psi_{01}}E_1\xrightarrow{\Psi_{12}}E_2$ over
$f_{01},f_{12}$, with $\Psi_{02}=\Psi_{12}\circ\Psi_{01}$ and
$f_{02}=f_{12}\circ f_{01}$,
\begin{equation}
A_{\Psi_{02}}(e;X)
=T^V\Psi_{12}\big|_{\Psi_{01}(e)}
\left(A_{\Psi_{01}}(e;X)\right)
+A_{\Psi_{12}}\left(\Psi_{01}(e);T_cf_{01}X\right).
\label{eq:pb-anomaly-composition}
\end{equation}
At the connection level this reads
$\mathfrak A_{\mathcal P_{02}}
=\mathcal P_{01}^*\mathfrak A_{\mathcal P_{12}}
+\mathrm d\kappa_{12}\circ\mathfrak A_{\mathcal P_{01}}$.
\status{ESTABLISHED}

\paragraph{Proof.}
Expand
$T\Psi_{01}(H^{\omega_0}_eX)
=H^{\omega_1}_{\Psi_{01}e}(Tf_{01}X)+A_{\Psi_{01}}(e;X)$
and apply $T\Psi_{12}$.  The horizontal summand contributes
$H^{\omega_2}_{\Psi_{02}e}(Tf_{02}X)+A_{\Psi_{12}}(\Psi_{01}e;Tf_{01}X)$ by
\eqref{eq:pb-coarse-horizontal-defect} at the second stage, and the vertical
summand contributes $T^V\Psi_{12}(A_{\Psi_{01}}(e;X))$ because a bundle
morphism preserves verticality.  Subtracting
$H^{\omega_2}(Tf_{02}X)$ gives \eqref{eq:pb-anomaly-composition}; the
connection-level form follows by
$T^V\Psi_{12}\circ\vartheta=\vartheta\circ\mathrm d\kappa_{12}$.  $\square$

The order in \eqref{eq:pb-anomaly-composition} is not cosmetic.  The earlier
defect takes values in $V\bar E_1$ and the later one in $VE_2$, so writing
$A_{\Psi_{02}}=A_{\Psi_{01}}+A_{\Psi_{12}}$ is a type error before any question
of correctness arises: the earlier defect must be pushed by the \emph{later}
vertical differential, and the later defect must be evaluated at the image point
on the \emph{pushed} base vector.  Correspondingly,
  $A_{\Psi_{02}}(e;X)=0$ holds if and only if
  \begin{equation}
  T^V\Psi_{12}\big|_{\Psi_{01}(e)}A_{\Psi_{01}}(e;X)
  =-A_{\Psi_{12}}\left(\Psi_{01}(e);T_cf_{01}X\right).
  \label{eq:pb-anomaly-composition-vanishing}
  \end{equation}
  Factorwise vanishing of $A_{\Psi_{01}}$ on $T\mathcal C_0$ and of
  $A_{\Psi_{12}}$ on $Tf_{01}(T\mathcal C_0)$ at the image points is sufficient,
  but it is not necessary because the two typed terms can cancel after the
  later vertical differential.  On the abelian three-level instance
  $f_{01}(x)=2x$, $f_{12}(y)=3y$ with local
connection forms $0$, $a_1dy$, and $a_2dz$ and identity fiber maps, the
  composition law reads $2a_1+2(3a_2-a_1)=6a_2$ identically in $(a_1,a_2)$;
  thus $a_2=0$ and $a_1\ne0$ give a vanishing composite defect while both
  factor defects are nonzero and cancel.  Dropping the base pushforward gives
  $a_1+3a_2$, and dropping the vertical differential fails likewise.
  \status{ESTABLISHED}

\begin{figure}[htbp]
\centering
\begin{tikzpicture}[>=stealth]
  \node[draw,rounded corners,fill=lightblue,minimum width=2.7cm,minimum height=0.9cm]
    (tc) at (0,2.1) {$T_c\mathcal C$};
  \node[draw,rounded corners,fill=lightgreen,minimum width=3.3cm,minimum height=0.9cm]
    (ve) at (5.4,2.1) {$V_{s(c)}E$};
  \node[draw,rounded corners,fill=lightblue,minimum width=2.7cm,minimum height=0.9cm]
    (tcb) at (0,0) {$T_{f(c)}\bar{\mathcal C}$};
  \node[draw,rounded corners,fill=lightgreen,minimum width=3.3cm,minimum height=0.9cm]
    (veb) at (5.4,0) {$V_{\bar s(f(c))}\bar E$};
  \draw[->,very thick,primaryblue] (tc)--node[above] {$D^\omega s$} (ve);
  \draw[->,very thick,primaryblue] (tcb)--node[below] {$D^{\bar\omega}\bar s$} (veb);
  \draw[->,very thick,sciencegreen] (tc)--node[left] {$T_cf$} (tcb);
  \draw[->,very thick,sciencegreen] (ve)--node[right] {$T^V\Psi$} (veb);
  \node[align=center,text width=7.0cm] at (2.7,-1.05)
    {$\mathcal D\Psi=0$: the first-jet square commutes};
\end{tikzpicture}
\caption{Naturality of the connection-split first jet.  The square commutes
exactly under the isotropy criterion \eqref{eq:pb-isotropy-criterion}.  When
the horizontal defect is nonzero, \eqref{eq:pb-covariant-jet-chain-rule}
records the missing vertical term and
\eqref{eq:pb-signed-base-comparison} records the two retained base tensors
$-\mathcal X_\Psi$ and $-\mathcal Q_\Psi$ rather than suppressing them.}
\label{fig:pb-pullback-naturality}
\end{figure}

\corollaryheading{Null directions under a parallel coarse morphism}{cor:pb-coarse-null-map}
If $\mathcal D\Psi=0$ along $s$, then
\begin{equation}
Tf(\ker D^\omega s)\subseteq\ker D^{\bar\omega}\bar s.
\label{eq:pb-coarse-null-inclusion}
\end{equation}
Thus a coarse morphism can create additional null directions; first-jet
naturality does not imply rank preservation.  \status{ESTABLISHED}

\paragraph{Proof.}
Set $D^\omega sX=0$ in
\eqref{eq:pb-covariant-jet-chain-rule} and use
$\mathcal D\Psi=0$.  $\square$

\section{Fisher defects and their composition law}
\label{sec:pb-fisher-defect}

Suppose now that the fiber map underlying $\Psi$ is the pushforward of a
parameter-independent Markov kernel $N$ between the two regular statistical
experiments, normalized in the strong sense that $N(x,\bar{\mathsf K})=1$ for
every $x$ rather than for almost every $x$, and that it is equivariant under
the represented gauge actions.  Three further hypotheses are consumed here and
none of them follows from affinity of the induced law map on the cone of
measures.  First, \emph{family closure}: the induced law map $N_\star$ must
satisfy $N_\star(\mathcal B)\subseteq\bar{\mathcal B}$, so that
$q:=N_\star|_{\mathcal B}$ is defined at all.  Second, \emph{smoothness of
$q$} between the two declared parametrized-measure models, together with a
jointly measurable and parameter-smooth version selection for
$p\mapsto T^V_p\Psi$.  Third, \Cref{hyp:pb-regular-models} is instantiated at
the coarse scale as well as the fine one; its clause that the represented
action is induced by a parameter-independent bimeasurable sample
re-coordinatization preserving the model is exactly the
$\widehat{\bar\rho}(\bar G)$-invariance of $\bar g^F$, and that invariance is
what cancels the context-dependent factor $\widehat{\bar\rho}(\varsigma(c))$ in
the local representative $\psi_c=\widehat{\bar\rho}(\varsigma(c))\circ q$ of
$\Psi$.  Only in a frame pair of the form $(u(c),\mathcal P(u(c)))$ is
$T^V\Psi$ literally $Tq$; in an arbitrary frame pair the identification
  requires that cancellation.  Finally, for the DQM transfer used below, assume
  a common $\sigma$-finite dominating measure $\mu$ for the selected fine
  family, with a fixed density version $p_\theta$ in the parameter neighborhood
  under discussion.  This common-domination tier is sufficient for the argument
  below; no broader nondominated transfer theorem is asserted here.  Jointly
  measurable, parameter-smooth conditional-score versions needed for pointwise
   bundle statements remain separate hypotheses. \status{HYPOTHESIS}

The coarse Fisher metric $\bar g^F$ can then be pulled back by $T^V\Psi$. Define
the vertical Fisher defect by
\begin{equation}
\Delta_F^\Psi
=g^F-(T^V\Psi)^*\bar g^F.
\label{eq:pb-vertical-fisher-defect}
\end{equation}
\status{DEFINITION}

\theoremheading{Fisher defect pullback and contraction}{thm:pb-pullback-fisher-defect}
Assume the preceding Markov hypotheses and
$\mathcal D\Psi=0$ along $s$.  Then
\begin{align}
\Delta_F^\Psi
&\succeq0,
\label{eq:pb-fisher-defect-positive}\\
h_s^\omega-f^*h_{\bar s}^{\bar\omega}
&=(D^\omega s)^*\Delta_F^\Psi\succeq0.
\label{eq:pb-pullback-fisher-contraction}
\end{align}
Equality on a base tangent vector $X$ holds exactly when the fine score in
the direction $D^\omega sX$ is measurable with respect to the coarse output.
\status{ESTABLISHED}

\paragraph{Proof.}
Work in a frame pair $(u(c),\mathcal P(u(c)))$, where $T^V\Psi=Tq$.  The step
that carries an arbitrary quadratic-mean path through the channel is not a
consequence of parameter independence by itself.  Let
$J_\theta(dx,dy)=P_\theta(dx)N(x,dy)$ and
$\nu(dx,dy)=\mu(dx)N(x,dy)$.  Then
$\mathrm dJ_\theta/\mathrm d\nu=p_\theta(x)$, so the fine DQM remainder has
exactly the same squared $L^2(\nu)$ norm in the joint experiment. Apply
Pollard's preservation theorem for measurable statistics \citep{Pollard2013}
to the deterministic statistic
$(x,y)\mapsto y$: the
pushed family $P_\theta N$ is DQM with directional score
$\bar\ell_w(y)=\E_{J_\theta}[\ell_w(X)\given Y=y]$.  The exact source and
hypothesis map are recorded in the Task~10 DQM-transfer source map.  Squared
Hellinger contraction is compatible corroboration, but it alone does not
establish this score identification.  Given the joint-experiment reduction
and Pollard's theorem, the law of total variance gives
\begin{equation}
\Delta_F^\Psi(u,u)
=\E\operatorname{Var}(\ell_u\given Y)\geq0,
\label{eq:pb-fisher-defect-score-variance}
\end{equation}
with equality exactly when $\ell_u$ is $Y$-measurable.  Under the
vanishing horizontal defect,
\eqref{eq:pb-covariant-tensor-naturality} applied to $\bar g^F$ gives
\[
f^*h_{\bar s}^{\bar\omega}
=(D^\omega s)^*(T^V\Psi)^*\bar g^F.
\]
Subtracting this identity from \eqref{eq:pb-fisher-pullback} proves
\eqref{eq:pb-pullback-fisher-contraction}.  $\square$

The guard $\mathcal D\Psi=0$ in \Cref{thm:pb-pullback-fisher-defect} may not be
dropped silently, and what survives when it is dropped is stated exactly rather
than gestured at.  Without the guard the vertical conclusion
\eqref{eq:pb-fisher-defect-positive} is unchanged, while the base conclusion
\eqref{eq:pb-pullback-fisher-contraction} is replaced by the exact signed
identity \eqref{eq:pb-signed-base-comparison} and its criterion
\eqref{eq:pb-positivity-criterion}; positivity then holds on exactly the
directions satisfying that criterion, and fails on the others.
\status{ESTABLISHED}

\corollaryheading{Coarse or meta-agent perceived geometry}{cor:pb-meta-perceived-geometry}
The coarse section $\bar s$ carries its own connection-relative semimetric
\begin{equation}
\bar h_{\bar s}^{\bar\omega}
=(D^{\bar\omega}\bar s)^*\bar g^F.
\label{eq:pb-meta-perceived-geometry}
\end{equation}
Under the exact section relation \eqref{eq:pb-coarse-related-sections} and
the compatibility condition $\mathcal D\Psi=0$ along $s$,
\begin{equation}
f^*\bar h_{\bar s}^{\bar\omega}
=(D^\omega s)^*(T^V\Psi)^*\bar g^F.
\label{eq:pb-meta-perceived-naturality}
\end{equation}
If the fiber map is also a genuine parameter-independent Markov channel,
then the fine-minus-coarse perceived-geometry defect is
\begin{equation}
h_s^\omega-f^*\bar h_{\bar s}^{\bar\omega}
=(D^\omega s)^*\Delta_F^\Psi\succeq0.
\label{eq:pb-meta-perceived-defect}
\end{equation}
\status{ESTABLISHED}

\paragraph{Proof.}
The first equation is the coarse-channel definition.  The second is
\eqref{eq:pb-covariant-tensor-naturality} with
$\bar\tau=\bar g^F$, and the third is
\eqref{eq:pb-pullback-fisher-contraction}.  $\square$

The inequality in \eqref{eq:pb-pullback-fisher-contraction} is
contravariant: the coarse tensor is pulled back to the fine tangent space.
No canonical pushforward of a Riemannian metric is being used.  If the coarse
operation is a restriction along an embedding, an energy precomposition, a
Galerkin map, or a parameter-dependent kernel, the Markov proof does not
apply.  \status{ESTABLISHED}

A perceived-geometry beta is not defined by subtracting tensors that live on
different level bases.  If the reference-space declarations of
\Cref{ch:renormalization} include base diffeomorphisms
$i_\ell:\mathcal C_\ell\to\mathcal C_\star$ and a block factor
$b_\ell>1$, then one may first define
\begin{equation}
\widetilde h_\ell=(i_\ell^{-1})^*h_\ell,
\qquad
\mathfrak B_{h,\ell}
=\frac{\widetilde h_{\ell+1}-\widetilde h_\ell}{\log b_\ell}
\quad\text{on }\mathcal C_\star.
\label{eq:pb-metric-beta-reference}
\end{equation}
This is the exact discrete perceived-geometry beta in that scheme.  On a
differentiable reference-space curve $s\mapsto\widetilde h_s$, its continuous
counterpart is $\beta_h(s)=\partial_s\widetilde h_s$. Both definitions are
conditional on the declared identifications and scale normalization.
\status{DEFINITION}

Without those identifications, $h_{\ell+1}-h_\ell$ is ill typed rather than a
beta function. \status{ESTABLISHED}

\theoremheading{Vertical Fisher defect cocycle}{thm:pb-fisher-defect-cocycle}
Let
$E_0\xrightarrow{\Psi_{01}}E_1\xrightarrow{\Psi_{12}}E_2$ be smooth bundle
morphisms over base maps $f_{01}:\mathcal C_0\to\mathcal C_1$ and
$f_{12}:\mathcal C_1\to\mathcal C_2$, carrying principal connections
$\omega_0,\omega_1,\omega_2$ and vertical Fisher metrics
$g_0^F,g_1^F,g_2^F$.  Whenever all three defects are defined,
\begin{equation}
\Delta_F^{\Psi_{12}\circ\Psi_{01}}
=\Delta_F^{\Psi_{01}}
+(T^V\Psi_{01})^*\Delta_F^{\Psi_{12}}.
\label{eq:pb-fisher-defect-cocycle}
\end{equation}
This identity is unconditional: it uses neither the connections nor any
section.  If both arrows are parameter-independent Markov coarse maps, the two
terms on the right are positive semidefinite.  \status{ESTABLISHED}

\paragraph{Proof.}
Contravariance of tensor pullback gives
\begin{align*}
\Delta_F^{\Psi_{12}\circ\Psi_{01}}
&=g_0^F-(T^V\Psi_{01})^*(T^V\Psi_{12})^*g_2^F\\
&=g_0^F-(T^V\Psi_{01})^*g_1^F
+(T^V\Psi_{01})^*
\left[g_1^F-(T^V\Psi_{12})^*g_2^F\right],
\end{align*}
which is \eqref{eq:pb-fisher-defect-cocycle}.  Positivity follows from
\eqref{eq:pb-fisher-defect-positive} for each Markov arrow.  $\square$

Pulling \eqref{eq:pb-fisher-defect-cocycle} back to the base does \emph{not}
give an additive identity in general, and the exact residual is available in
closed form.  Fix sections $s_0,s_1,s_2$ related at each stage, write
$\delta_{jk}=(D^{\omega_j}s_j)^*\Delta_F^{\Psi_{jk}}$ for the three base
defects, and along $s_0$ write
$v_X=T^V\Psi_{01}(D^{\omega_0}s_0X)$, $A_X=A_{\Psi_{01}}(s_0;X)$, and
$\bar u_X=D^{\omega_1}s_1(Tf_{01}X)=v_X+A_X$.

\theoremheading{Base Fisher defect cocycle: exact residual and sharp criterion}{thm:pb-base-defect-cocycle}
The residual
$\mathcal N:=\delta_{02}-\delta_{01}-f_{01}^*\delta_{12}$ has the three
equivalent exact forms
\begin{align}
\mathcal N(X,Y)
&=\Delta_F^{\Psi_{12}}(v_X,v_Y)-\Delta_F^{\Psi_{12}}(\bar u_X,\bar u_Y),
\label{eq:pb-base-cocycle-residual-a}\\
&=-\left[\Delta_F^{\Psi_{12}}(\bar u_X,A_Y)
+\Delta_F^{\Psi_{12}}(A_X,\bar u_Y)\right]
+\Delta_F^{\Psi_{12}}(A_X,A_Y),
\label{eq:pb-base-cocycle-residual-b}\\
&=-\left[\Delta_F^{\Psi_{12}}(v_X,A_Y)
+\Delta_F^{\Psi_{12}}(A_X,v_Y)\right]
-\Delta_F^{\Psi_{12}}(A_X,A_Y).
\label{eq:pb-base-cocycle-residual-c}
\end{align}
The sign of the quadratic term is coupled to the choice of argument in the
cross terms: the coarse jet $\bar u$ carries $+\Delta_F^{\Psi_{12}}(A,A)$ and
the pushed fine jet $v$ carries $-\Delta_F^{\Psi_{12}}(A,A)$.  Mixing the two
conventions inflates the residual by exactly
$2\Delta_F^{\Psi_{12}}(A_X,A_Y)$, which is two copies of the second-arrow
vertical Fisher defect evaluated on the first-arrow horizontal defect.
Consequently the sharp base cocycle
$\delta_{02}=\delta_{01}+f_{01}^*\delta_{12}$ holds \emph{if and only if}
\begin{equation}
\Delta_F^{\Psi_{12}}(v_X,v_X)
=\Delta_F^{\Psi_{12}}(\bar u_X,\bar u_X)
\quad\text{for every }X,
\label{eq:pb-base-cocycle-criterion}
\end{equation}
equivalently $\Delta_F^{\Psi_{12}}(A_X,2v_X+A_X)=0$ for every $X$.  If the
second arrow is a Markov coarse map, so that
$\Delta_F^{\Psi_{12}}\succeq0$, this is equivalently equality of the two
$\Delta_F^{\Psi_{12}}$-seminorms.
\status{ESTABLISHED}

\paragraph{Proof.}
Apply $(D^{\omega_0}s_0)^*$ to \eqref{eq:pb-fisher-defect-cocycle} to get
$\delta_{02}=\delta_{01}+\Delta_F^{\Psi_{12}}(v_X,v_Y)$, and note that by
definition $f_{01}^*\delta_{12}(X,Y)=\Delta_F^{\Psi_{12}}(\bar u_X,\bar u_Y)$.
Subtracting gives \eqref{eq:pb-base-cocycle-residual-a}.  Substituting
$\bar u=v+A$ in the second term gives
\eqref{eq:pb-base-cocycle-residual-c}, and substituting $v=\bar u-A$ in the
first gives \eqref{eq:pb-base-cocycle-residual-b}.  Since $\mathcal N$ is
symmetric, vanishing on the diagonal is vanishing, and
\eqref{eq:pb-base-cocycle-residual-a} on the diagonal is exactly
\eqref{eq:pb-base-cocycle-criterion}.  $\square$

Three conditions are sufficient for
\eqref{eq:pb-base-cocycle-criterion}, in decreasing strength: the stage-one
horizontal defect vanishes, which by \Cref{thm:pb-isotropy-criterion} is the
isotropy condition \eqref{eq:pb-isotropy-criterion} at the first arrow; the
second arrow is Fisher lossless, $\Delta_F^{\Psi_{12}}=0$; or, weakest and new,
$A_X\in\operatorname{rad}\Delta_F^{\Psi_{12}}$ for every $X$, so that the
stage-one anomaly is invisible to the stage-two information loss.  The third
suffices because $\mathcal N$ is linear in the product
$\Delta_F^{\Psi_{12}}A$.  None of the three is necessary, since
\eqref{eq:pb-base-cocycle-criterion} is an equality of quadratic forms and admits
accidental solutions: on the three-level normal location instance with fibers
of variance $1,2,3$, sections $x\mapsto x$, $y\mapsto y/2$, $z\mapsto z/6$,
base maps $f_{01}(x)=2x$ and $f_{12}(y)=3y$, and local connection forms $0$,
$a_1dy$, $a_2dz$, the residual is
$\mathcal N=-\tfrac23a_1(a_1+1)$, which vanishes at $a_1=-1$ where the
stage-one horizontal defect equals $-2\ne0$.  On the same data the
convention-mixed expression evaluates to $\tfrac23a_1(a_1-1)$, exceeding
$\mathcal N$ by $\tfrac43a_1^2=2\Delta_F^{\Psi_{12}}(A,A)$, and at
$a_1=\tfrac1{10}$ the two numbers are $-\tfrac{11}{150}$ and $-\tfrac3{50}$.
\status{ESTABLISHED}

Two distinct hypotheses are involved here and conflating them is the reading
this chapter forbids.  The sharp base cocycle needs only the stage-one
condition \eqref{eq:pb-base-cocycle-criterion}.  Reading the middle term
$\delta_{12}$ as the stage-two information loss $h_1-f_{12}^*h_2$ requires in
addition that the stage-two horizontal defect vanish, by
\Cref{thm:pb-signed-base-comparison} applied at the second arrow.  Only when
both hold does \eqref{eq:pb-fisher-defect-cocycle} become an additive
decomposition of information loss on the base.  \status{ESTABLISHED}

The Amari defect
\begin{equation}
\Delta_{\mathcal T}^\Psi
=\mathcal T-(T^V\Psi)^*\bar{\mathcal T},
\label{eq:pb-amari-defect}
\end{equation}
obeys the same algebraic cocycle identity.  It has no analogous
positive-semidefinite order because a nonzero cubic changes sign under
$u\mapsto-u$.  Exact Amari naturality therefore requires a
$\mathcal T$-preserving statistical morphism; Markov coarse-graining alone
does not supply a monotonicity statement for the cubic tensor.
\status{ESTABLISHED}

\section{Boundary of the construction}
\label{sec:pb-boundary}

The chapter has constructed vertical tensors and their functorial defects.  A
scalar gauged sigma energy would additionally require a base cometric, a base
density, channel weights, boundary conditions, and a decision about whether
the connection is fixed or dynamical.  None is selected by
$h_s^\omega$ or $c_s^\omega$.  \status{NOT-CLAIMED}

Likewise, vertical Fisher length is invariant under reparameterization of a
fixed regular curve, but it does not generate an orbit, identify a physical
time coordinate, or compare independently evolved fine and coarse paths.
Comparing such paths would require a semiconjugacy of their vector fields in
addition to the pointwise contraction
\eqref{eq:pb-pullback-fisher-contraction}.  That interface is developed in
\Cref{ch:relational-inference}, where the semiconjugacy criterion, its
sufficient natural-gradient conditions, and the exact duration relation are
proved on a declared configuration tier; what remains open for the specific
renormalization maps of this work is the obligation recorded at
\Cref{prop:hist-natural-gradient-sufficiency}, namely exhibiting the functional
compatibility and horizontal conformality of the declared coarse configuration
map.  \status{OPEN}

The scale index carried by \eqref{eq:pb-fisher-defect-cocycle} and by
\Cref{thm:pb-base-defect-cocycle} is renormalization depth.  It is not an
inference-orbit coordinate and not a duration, and no result of this chapter
converts it into either.  \status{NOT-CLAIMED}

The general coarse-graining chapter uses the typed distinction exposed here.
Its normalized Markov channels inherit the Fisher contraction theorem.  Its energy
precompositions and operator restrictions require their own probability and
closure arguments and do not inherit that theorem by terminology.
\status{ESTABLISHED}
